\section{Materials and methods}
\label{sec:methods}

\subsection{Datasets}
\label{sec:datasets}

The primary benchmark is BP3C50ID, the non-redundant split distributed with
BepiPred-3.0, in which training and external test antigens are separated at a
50\% sequence identity threshold. The external test split contains \BepipredNumAntigens{} antigens,
\BepipredResidues{} residues and \BepipredPositives{} positive residues, an
epitope-residue rate just under one residue in five. All BP3C50ID results below
are on this independent benchmark; no model or policy parameter is fitted on
it.

A second evaluation uses a subset derived from the Immune Epitope Database
\citep{vita2024iedb}, holding \IedbSubsetTest{} test antigens
(\NestedIedbResidues{} residues, \NestedIedbPositives{} positive) alongside
\IedbSubsetTrain{} training antigens. The subset is drawn from the
\texttt{\IedbSourceFile{}} collection distributed alongside BP3C50ID, which
holds \IedbSourceSequences{} IEDB sequences selected not to share more than
20\% sequence identity with the BP3 data, and is then restricted to sequences of
\IedbLengthMin{}--\IedbLengthMax{} residues so that the full policy grid remains
runnable on a workstation. Every sequence in the subset is present verbatim in
that source collection.

Its purpose is to test whether conclusions drawn from \BepipredNumAntigens{}
antigens survive on a different and larger collection, which requires the two
to be genuinely independent. Independence holds by construction through the
identity filter. We also tested it directly against all
\IedbBpThreeCompared{} BP3C50ID sequences and found \IedbIdOverlap{} shared
identifiers, \IedbExactOverlap{} shared sequences, and a maximum
\IedbKmerSize{}-mer containment of \IedbMaxContainment{}. No antigen in the
subset therefore shares even a single \IedbKmerSize{}-residue run with any
BP3C50ID antigen.

\subsection{Predictors}
\label{sec:predictors}

We evaluate two families of local and remote predictor of conformational B-cell
epitopes, and hold the distinction between them fixed across every policy.

The \emph{heuristic family} uses handcrafted per-residue sequence-context
features. Its local member is a lightweight residue classifier; its remote
member is an extremely randomised trees refiner that also blends the
local score and applies sequence smoothing, with the blend weight and smoothing
parameters chosen on a validation split.

The \emph{PLM family} replaces the features with frozen embeddings from ESM-2
\citep{lin2023esm}, specifically the \texttt{esm2\_t6\_8M\_UR50D}
checkpoint. The local member places a logistic residue head on those embeddings;
the remote member adds contact-derived features and the same trees refiner.
A region-constrained variant post-processes the remote probabilities into
contiguous regions using contact density, with the region parameters selected on
validation data. We use the 8M-parameter checkpoint because it is what runs on
the hardware available for this study. That checkpoint is the smallest of the released
ESM-2 models. The PLM results here should therefore be read as a lower bound on what
the family can do, not as a characterisation of language-model-based epitope
prediction in general.

\subsection{Routing policies}
\label{sec:policies}

A routing policy decides, per antigen, whether the remote model is invoked.
Five are compared in each family.

\begin{description}
  \item[local-only] The local model answers every antigen. No sequence is
  transmitted.
  \item[edge-only] Every antigen is sent to the remote model.
  \item[adaptive-split] The local model runs first; if its mean per-residue
  uncertainty exceeds a threshold, the antigen is re-processed remotely and the
  remote prediction is used. The threshold and a routing budget are selected on
  the validation split from a grid of \ThresholdGridSize{} candidate values,
  subject to a maximum routed fraction of \RouteBudgetMax{}.
  \item[carbon-guard] As adaptive-split, with an additional constraint that
  suppresses routing when the modelled carbon cost of the remote call is not
  justified by the expected quality gain.
  \item[conformal-guard (PLM family)] As adaptive-split, with a conformal
  criterion \citep{angelopoulos2021gentle} replacing the raw uncertainty
  threshold.
\end{description}

Every selectable parameter in this section is chosen on a validation split
carved out of training data at a fraction of \ValidationFraction{}, grouped by
sequence so that no antigen contributes residues to both sides. No routing
parameter sees the external test split.

\subsection{Cost accounting}
\label{sec:cost-accounting}

Four cost quantities are recorded per antigen, and they differ in how directly
they are observed; Table~\ref{tab:provenance} states this for each.

Latency is measured, as wall-clock model-prediction time on the executing
device. Bytes transferred are computed from the raw sequence length. Raw-sequence
share rate is a count: the fraction of test antigens whose sequence is
transmitted at all, which is the quantity a laboratory worried about disclosure
would actually audit.

Energy and carbon are modelled, not metered. Operational energy is the product
of a fixed device power figure and the measured latency; operational carbon
converts that with a grid-intensity constant of
\SI{708}{\gram\per\kilo\watt\hour}; embodied carbon is a per-request
amortisation of a device lifetime figure. The local device is modelled at
\SI{34}{\watt} and the remote at \SI{145}{\watt}.

The embodied term also carries a device-utilisation fraction. We apply it
so that a request is charged only the share of the device's lifetime footprint
it actually occupies: \CarbonLocalFraction{} locally and \CarbonEdgeFraction{}
remotely. The two fractions reduce the per-request embodied charge from
\EmbodiedDeclaredLocal{} to \EmbodiedEffectiveLocal{}~gCO\textsubscript{2}e
locally and from \EmbodiedDeclaredEdge{} to \EmbodiedEffectiveEdge{} remotely.
We state them explicitly because they are the single largest determinant of
every carbon figure in this paper. As Section~\ref{sec:limitations}
discusses, the comparison between policies is not fully insensitive to them.

The balance between those two terms matters for reading the results, so we
state it plainly. At the latencies observed here the operational term is a
very small part of each total: a local pass contributes about
\OperationalLocal{}~gCO\textsubscript{2}e and a remote pass about
\OperationalEdge{}, against
per-request totals near \LocalOnlyCarbon{} and \EdgeOnlyCarbon{}. The modelled
carbon of a policy is therefore governed almost entirely by how many
device-requests it charges, through the amortised embodied term, and only
marginally by how long those requests take. Carbon differences between
policies here are differences in request counts, not in compute. We make no
claim about metered power anywhere in this paper.

\subsection{Evaluation protocol for routing policies}
\label{sec:eval-protocol}

The full pipeline (validation split, parameter selection, training and
external-split evaluation) is repeated for \NumSeeds{} seeds
(\SeedFirst{}--\SeedLast{}).

Repeating a pipeline over seeds on a fixed test split answers a narrow
question: how much a number moves when the same antigens are scored again. It
says nothing about how the number would move on different antigens, which is
what a reader wants to know. Our primary analysis therefore resamples both.
For each of \AntigenBootstrapRepeats{} replicates we draw \AntigenCount{}
antigens with replacement and \NumSeeds{} seeds with replacement, recompute
every policy's pooled residue-level metrics on that resampled set. We take
the 2.5th and 97.5th percentiles of the resulting distribution as a 95\%
interval. One antigen draw is shared across all policies within a replicate,
so paired differences stay paired. Intervals are percentile intervals, with no
bias correction or studentisation, and the resampling is driven by a fixed
seed recorded in the released artifact.

We report seed-only intervals alongside these in Section~\ref{sec:results}
only to show how different the two are; no claim in this paper rests on them,
and the exact sign test over seeds that an earlier version of this work had relied
on is not used.

Two cautions apply to how we read these intervals. First, with
\AntigenCount{} antigens a percentile cluster bootstrap is known to
under-cover for non-linear statistics such as MCC and ROC-AUC, so the
intervals should be read as indicative rather than exact. Second, several of
the quantities we report are not random in sign. A policy that routes only
some antigens cannot be slower than one that routes all of them, and cannot
transmit more sequences. For latency and disclosure the interval therefore describes
the magnitude of a difference whose direction is fixed by construction.
Reporting that such an interval excludes zero would assert nothing. We say so
where it applies, and we reserve the language of detection for the quality
metrics, where the sign is genuinely in question. There we treat an interval
containing zero as no detected difference, and say so rather than describing
such a comparison as a trend.

\subsection{Two protocols for selecting a fusion weight}
\label{sec:fusion-protocol}

Score fusion combines the reference predictor's per-residue score with ours
(CARBON-BCE-R, the local-only predictor) into a meta-predictor,
a convex blend governed by a mixing weight \(w\). The two predictors are calibrated
differently, so a blend of their raw outputs would be dominated by whichever
has the wider spread. We therefore map each onto \([0,1]\) through its own
empirical distribution function before blending. We estimate that distribution
on the training folds and apply it to the held-out fold, like any other fitted
transform. The mapping is not a rank transform computed over the whole split,
which would let held-out residues influence their own normalisation. One consequence is worth
stating, because it is easy to assume otherwise. A transform fitted on training
data is not exactly
rank-preserving on held-out data, so a fold's ROC-AUC after mapping can differ
slightly from the same predictor's raw ROC-AUC. The effect is small, under
\(10^{-3}\) per fold in our runs, but it is not zero. For this reason we report
per-fold paired differences rather than pooled ranking metrics. A decision
threshold then converts
blended scores into calls. Both \(w\) and the threshold are tuning parameters.
The whole question is where they come from.

\paragraph{Notation.} Let \(\mathcal{A}\) be the set of \(n\) test antigens,
\(R_a\) the residues of antigen \(a\), and \(y_{a,i}\in\{0,1\}\) the epitope
label of residue \(i\in R_a\). Each residue carries two raw scores:
\(r_{a,i}\) from the reference predictor and \(c_{a,i}\) from ours. For a set
of antigens \(T\subseteq\mathcal{A}\), let \(\hat{F}^{\mathrm{ref}}_{T}\) and
\(\hat{F}^{\mathrm{cand}}_{T}\) be the empirical distribution functions of the
two raw score sets over the residues of \(T\) alone. The blended score of
residue \(i\) of antigen \(a\) at weight \(w\), normalised through \(T\), is
\begin{equation}
\label{eq:blend}
s_{a,i}(w,T) \;=\; w\,\hat{F}^{\mathrm{ref}}_{T}\!\bigl(r_{a,i}\bigr)
\;+\;(1-w)\,\hat{F}^{\mathrm{cand}}_{T}\!\bigl(c_{a,i}\bigr),
\qquad w\in\mathcal{W}\subset[0,1],
\end{equation}
so that \(w\) is the weight on the \emph{reference} predictor and \(1-w\) the
weight on ours; \(\mathcal{W}\) is the finite grid of candidate weights that is
scanned. We write \(s_{T}(w,U)\) for the vector of \(s_{a,i}(w,U)\) over the
residues of all antigens in \(T\), and \(y_{T}\) for the corresponding labels.

The threshold is not a free scalar either. For a score vector \(s\) and a
target positive-call rate \(\rho\),
\begin{equation}
\label{eq:rate-threshold}
\tau(s,\rho)\;=\;s_{(k)},\qquad k=\operatorname{round}\!\bigl(\rho\,\lvert s\rvert\bigr),
\end{equation}
the \(k\)-th largest entry of \(s\): the operating point is a constrained
quantile, chosen so the blend emits the same number of positive calls as the
reference. We write \(\rho_{T}\) for the positive-call rate the reference
predictor produces at its own published threshold on the residues of \(T\), and
\(\widehat{M}(y,s,\tau)\) for the empirical selection metric of scores \(s\) against labels \(y\)
at operating point \(\tau\). The two protocols below use exactly these objects
and differ only in which data the maximisation in \(w\), the distribution
functions, and the threshold are allowed to see.

\paragraph{Test-derived selection.} The protocol used in the earlier version of
this work optimizes both tuning parameters on the evaluation split, reporting
\begin{equation}
\label{eq:test-derived}
\widehat{\Delta}_{\mathrm{test}}
=\max_{w\in\mathcal{W}}
\widehat{M}\!\Bigl(y_{\mathcal{A}},\,s_{\mathcal{A}}(w,\mathcal{A}),\,
\tau\bigl(s_{\mathcal{A}}(w,\mathcal{A}),\rho_{\mathcal{A}}\bigr)\Bigr)
\;-\;
\widehat{M}\!\Bigl(y_{\mathcal{A}},\,\hat{F}^{\mathrm{ref}}_{\mathcal{A}}(r),\,
\tau\bigl(\hat{F}^{\mathrm{ref}}_{\mathcal{A}}(r),\rho_{\mathcal{A}}\bigr)\Bigr).
\end{equation}
Every argument on the right-hand side is the whole evaluation split: the grid
maximum, the normalisation and the threshold all see \(y_{\mathcal{A}}\). Both
tuning parameters are thus functions of the test labels, and the margin the
protocol reports over the stronger single predictor is an over-optimistic
estimate.

\paragraph{Nested selection.} The corrected protocol is nested
cross-validation with the antigen as the unit of resampling. The outer loop
leaves out one antigen at a time, so that \(T_a=\mathcal{A}\setminus\{a\}\)
carries every quantity fitted for fold \(a\):
\begin{align}
w_a &= \operatorname*{arg\,max}_{w\in\mathcal{W}}
 \widehat{M}\!\Bigl(y_{T_a},\,s_{T_a}(w,T_a),\,
 \tau\bigl(s_{T_a}(w,T_a),\rho_{T_a}\bigr)\Bigr), \label{eq:inner-selection}\\
\tau_a &= \tau\bigl(s_{T_a}(w_a,T_a),\,\rho_{T_a}\bigr). \label{eq:fold-threshold}
\end{align}
The inner loop \eqref{eq:inner-selection} scores the same parameter grid on the
training folds of that fold; it is a grid scan over those folds, evaluated on
the very residues it is fitted on, not a further resampling split inside them.
The weight and the threshold \eqref{eq:fold-threshold} are then applied,
unchanged, to the held-out antigen, and the outer estimate is the mean of the
per-antigen paired differences,
\begin{equation}
\label{eq:nested}
\widehat{\Delta}_{\mathrm{nested}}
=\frac{1}{n}\sum_{a\in\mathcal{A}}
\Bigl[\,
\widehat{M}\!\bigl(y_a,\,s_a(w_a,T_a),\,\tau_a\bigr)
-\widehat{M}\!\Bigl(y_a,\,\hat{F}^{\mathrm{ref}}_{T_a}(r_a),\,
\tau\bigl(\hat{F}^{\mathrm{ref}}_{T_a}(r_{T_a}),\rho_{T_a}\bigr)\Bigr)
\Bigr].
\end{equation}
Comparing \eqref{eq:nested} with \eqref{eq:test-derived}: the metric, the
blend and the threshold rule \eqref{eq:rate-threshold} are the same functions; only the subscript on the
data they are evaluated with has changed. Nothing derived from the
held-out antigen enters the configuration applied to it: not its labels, not its
score distribution, not its own quantiles. Independence of that kind is what makes the
resulting predictor deployable, since it needs no knowledge of a new antigen
before scoring its first residue. We report
\(\widehat{\Delta}_{\mathrm{nested}}\) with a 95\% bootstrap interval over
antigens, a Wilcoxon signed-rank test and an exact sign test. We also report
the weight \(w_a\) chosen in every fold, not its median alone.

\paragraph{What nested selection does and does not buy.} The property that
holds is an independence one. \emph{Claim:} given the data of \(T_a\), the
out-of-fold scores and calls for antigen \(a\) in \eqref{eq:nested} are fixed
functions of that antigen's own inputs, and in particular do not depend on
\(y_a\). \emph{Proof:} by \eqref{eq:inner-selection} and
\eqref{eq:fold-threshold}, \(w_a\), \(\tau_a\),
\(\hat{F}^{\mathrm{ref}}_{T_a}\) and \(\hat{F}^{\mathrm{cand}}_{T_a}\) are
functions of \((r,c,y)\) restricted to \(T_a\); substituting them into
\eqref{eq:blend} makes \(s_{a,i}(w_a,T_a)\) a function of \((r_{a,i},c_{a,i})\)
and those four quantities only, and the call is
\(\mathbf{1}[s_{a,i}(w_a,T_a)\ge\tau_a]\). So \(y_a\) enters \eqref{eq:nested}
only through \(M\), as the labels that already-fixed predictions are scored
against. The single assumption is the one the design enforces: folds are
defined by antigen, so no residue of \(a\) belongs to \(T_a\).

What that licenses is reading \eqref{eq:nested} as an estimate of the
performance of the selection \emph{procedure} on an antigen it has not seen.
What it does not license is calling \(\widehat{\Delta}_{\mathrm{nested}}\) an
unbiased estimate of that performance, and we do not claim it is, for two
reasons. First, the inner selection \eqref{eq:inner-selection} is
resubstitution: the metric is maximised on the same training residues at the
threshold fitted on those residues, so \(w_a\) is itself the argmax of a noisy
in-sample score, and \eqref{eq:nested} inherits whatever optimism that choice
carries into the held-out fold. Second, any two folds share \(n-2\) antigens,
so the per-antigen differences in \eqref{eq:nested} are not independent; the
spread of the leave-one-out mean is smaller than independent folds would give,
which is why the percentile interval over folds under-covers rather than
attaining its nominal level. The nested number is the better of the two
estimates on offer, not a guaranteed-unbiased one.

Both protocols are implemented in the same module against the same loaders and
metric code, so that a difference between their outputs is attributable to the
protocol rather than to an implementation difference. Section~\ref{sec:replication}
verifies this by reproducing the earlier figures under the test-derived
protocol, to within \TunedBpThreeReplicationMaxDiff{}.
